\SourceLeaf{270}%
Consequently, using 5 to eliminate either \(\mathrm{d}U\) or \(\mathrm{d}w\), we have
\[
\frac{\mathrm{d}\log m}{\mathrm{d}w}=\frac{(1-ee)(\alpha\cos U-\cos w)}{(1-ee\cos w^2)\sin w} \quad .\ .\ .\ .\ .\ .\ .\ .\ .\ .\ .\ .\quad(6)
\]
\[
\frac{\mathrm{d}\log m}{\mathrm{d}U}=\operatorname{cotg}U-\frac{\cos w}{\alpha\sin U}=\frac{\alpha\cos U-\cos w}{\alpha\sin U} \quad .\ .\ .\ .\ .\ .\ .\quad(7)
\]

Differentiating equation 7 once more gives
\[
\begin{aligned}
\frac{\mathrm{dd}\log m}{\mathrm{d}U^2}&=-\frac{1}{\sin U^2}+\frac{\cos U\cos w}{\alpha\sin U^2}+\frac{\sin w}{\alpha\sin U}\mathord{.}\frac{\mathrm{d}w}{\mathrm{d}U}\\
&=-\frac{1}{\sin U^2}+\frac{\cos U\cos w}{\alpha\sin U^2}+\frac{(1-ee\cos w^2)\sin w^2}{\alpha\alpha(1-ee)\sin U^2} \quad .\ .\ .\ .\ .\ .\quad(8)
\end{aligned}
\]

Suppose that \(m\) is to equal unity at a specified latitude (standard latitude) and to differ from 1 at other latitudes only by quantities of the third order, with the differences in latitude treated as quantities of the first order. Let \(P\) denote the standard latitude on the ellipsoid and \(Q\) the corresponding latitude on the sphere. Then for \(w=90^0-P\), \(U=90^0-Q\), equations 4, 7 and 8 require the following.
\[
A=\frac{a\cos P}{\alpha\cos Q\BareRoot{(1-ee\sin P^2)}} \quad .\ .\ .\ .\ .\ .\ .\ .\ .\ .\ .\ .\ .\ .\quad(9)
\]
\[
\alpha\sin Q=\sin P \quad .\ .\ .\ .\ .\ .\ .\ .\ .\ .\ .\ .\ .\ .\ .\ .\ .\ .\ .\quad(10)
\]
\[
0=1-\frac{\sin P\sin Q}{\alpha}-\frac{(1-ee\sin P^2)\cos P^2}{\alpha\alpha(1-ee)}
\]
Or, substituting in the last equation the value of \(\sin Q\) from 10, we obtain
\[
\alpha\alpha=1+\frac{ee\cos P^4}{1-ee} \quad .\ .\ .\ .\ .\ .\ .\ .\ .\ .\ .\ .\ .\ .\ .\ .\quad(11)
\]

This equation thus gives \(\alpha\) once a particular value of \(P\) has been chosen. Equation 10 can then be used to determine \(Q\), and equation 9 to determine \(A\). Finally, \(k\) is obtained by substituting \(w=90^0-P\), \(U=90^0-Q\) into the general equation 3, giving
\[
k=\frac{\operatorname{tang}(45^0+\frac{1}{2}P)^\alpha}{\operatorname{tang}(45^0+\frac{1}{2}Q)}\mathord{.}(\tfrac{1-e\sin P}{1+e\sin P})^{\frac{1}{2}\alpha e} \quad .\ .\ .\ .\ .\ .\ .\ .\quad(12)
\]

\Article{4}
\indent One will want to calculate the constants \(A,\alpha,k\) and the standard latitude on the sphere \(Q\) from \(P\) and \(e\) with particular care and precision, since all these quantities are to be regarded as the basis for
